\documentclass{article}

\usepackage[preprint]{neurips_2024}

\usepackage[utf8]{inputenc}
\usepackage[T1]{fontenc}
\usepackage{hyperref}
\usepackage{url}
\usepackage{booktabs}
\usepackage{amsfonts}
\usepackage{amsmath}
\usepackage{amssymb}
\usepackage{microtype}
\usepackage{xcolor}
\usepackage{graphicx}
\usepackage{multirow}
\usepackage{makecell}
\setcitestyle{numbers,square}

\newcommand{\todo}[1]{\textcolor{red}{\textbf{[TODO: #1]}}}
\newcommand{\Mc}{M_{\mathrm{c}}}
\newcommand{\Ct}{C_t}
\newcommand{\mt}{m_t}
\newcommand{\Min}{M_{\mathrm{in}}}
\newcommand{\Mj}{M_{\mathrm{judge}}}
\newcommand{\Mnew}{M_{\mathrm{new}}}
\newcommand{\dnm}{\Delta_{\mathrm{nm-m}}}
\newcommand{\dsh}{\Delta_{\mathrm{sh-m}}}
\newcommand{\CEmem}{\mathrm{CE}_{\mathrm{mem}}}
\newcommand{\CEnomem}{\mathrm{CE}_{\mathrm{nomem}}}
\newcommand{\CEsh}{\mathrm{CE}_{\mathrm{shuffle}}}
\newcommand{\CEfloor}{\mathrm{CE}_{\mathrm{floor}}}
\newcommand{\elift}{\Delta_{\mathrm{evidence}}}
\newcommand{\amem}{\alpha_{\mathrm{mem}}}

\title{When Does Recurrent Memory Help?\\Six Failure Modes and a Methodology for Honest Evaluation}

\author{%
  Anonymous Authors\\
  \texttt{anonymous@example.com}
}

\begin{document}

\maketitle

\begin{abstract}
We began with a simple question: why did a recurrent memory module
sometimes look useful, sometimes invisible, and sometimes actively
harmful on the same synthetic callback task?  In a 16-version campaign
on Memory Residuals, a slot-attention memory module attached to
Qwen3-0.6B and Qwen3-1.7B, the answer was not a single bug.  It was an
evaluation problem.  Token-averaged chain evaluation diluted a
localized callback gain by roughly two orders of magnitude; frozen and
joint training measured different causal pathways; the D4v2 callback
template leaked a 32-way category prior; more than half of the D4v2
callback windows exposed evidence directly to the LM path; evidence
redaction had to be audited before a zero lift could be trusted; and a
random-ID corpus removed the prior but exposed a writer objective that
still did not force chain-specific bindings into $\Mc$.  We turn these
failures into a methodology: callback-aware evaluation, evidence-
redacted memory floors, empirical and base-model prior audits,
window-leakage counts, and a test-time readout audit that asks whether
the writer encoded the answer at all.  The result is deliberately
negative but actionable: most apparent memory gains on D4v2 are
shortcuts, while D5 shows that removing shortcuts is not sufficient
without writer-side extractive supervision.  We close with an F2
prescription: NIAH-style writer warmup that gives $\Mc$ a direct
chain-specific training signal before the LM head can learn a shortcut.
\end{abstract}

\section{Introduction}

Recurrent memory is attractive because it promises a bounded-cost state
that survives across sessions.  In a conversational assistant, that
state should remember a user's code, preference, or updated fact even
when the evidence is no longer in the language model's local attention
window.  The hard question is not whether an augmented model can reduce
loss somewhere.  It is whether the recurrent state is the causal path
that carries the chain-specific binding.

This paper is an empirical post-mortem of that question.  We study
Memory Residuals: a fixed-size memory state $\Mc$ written once per
session and read back into a Qwen3 backbone through a depth-routed
attention residual.  Across versions v11--v16, the system produced
encouraging numbers, contradictory numbers, and finally clean negative
numbers.  The campaign changed our interpretation of the result.  The
main contribution is not a state-of-the-art memory model; it is a set of
failure modes and diagnostics that determine when a memory result is
honest.

\paragraph{The Running Task.}
Most of the campaign uses synthetic multi-session chains.  An evidence
session states a binding; later, a callback session asks for it.  The
model is scored on the answer tokens.  In D4v2 the binding is a
closed-set persona fact such as ``favorite color.''  In D5 the answer is
a per-chain random alphanumeric code and the evidence-session LM loss is
masked so the callback span is the only supervised answer location.
The intended invariant is simple: if the evidence is outside the LM
window and the callback does not reveal the answer, only $\Mc$ should
carry the value.

\paragraph{What Went Wrong.}
Five concrete findings forced a reframe.  First, the standard
\texttt{eval\_chain.py} metric averages over full score windows.  On
v14k, callback-aware evaluation reports $\dnm^{\mathrm{cb}}=+1.44$ nats,
but standard full-window evaluation reports $\dnm \approx 0$ because
the callback answer tokens are only about 1.4\% of the scored window.
Second, the frozen-backbone and joint-training regimes measure
different causal systems: the frozen 0.6B path can show apparent memory
help, while the joint 0.6B path learns the D4v2 prior directly and
collapses $\elift$ to zero.  Third, D4v2 leaks a category prior: the
callback ``Your favorite COLOR is \_\_'' reduces the answer to a
32-item closed set, and the empirical prior has CE $2.91$ nats per
callback-mask token.  v15b reaches $2.97$ nats, within $0.06$ nats of
that prior.  Fourth, D5 removes that closed-set cue but does not cure
the writer objective: v16a still has near-zero $\elift$ because the
writer is not forced to encode the chain-specific code.  Fifth, the
evidence-redacted floor is clean: Audit A2 shows default, cross-chain,
skip, and zero redactions agree within $0.05$ nats, so a zero lift is
not an artifact of leaky redaction.

\paragraph{Contributions.}
We contribute four evaluation tools and one prescription.
\begin{enumerate}
\item \textbf{Callback-aware standalone evaluation.}  \texttt{tools/eval\_callback.py}
scores only annotated callback answer tokens and reports an evidence-
redacted $\Mc$ floor.
\item \textbf{Multi-axis leakage audits.}  We measure empirical
$P(\mathrm{item}\mid\mathrm{category})$ CE, frozen-base callback CE,
window-leakage rates, and redaction validity.
\item \textbf{D5 random-ID stress test.}  The D5 corpus removes the D4v2
category prior and masks evidence-session LM loss, exposing whether the
writer objective alone stores the binding.
\item \textbf{TTT-readout audit.}  A planned diagnostic freezes the
model and optimizes only readout-side parameters on callback tokens to
test whether the writer encoded the answer in $\Mc$.
\item \textbf{F2 prescription.}  We propose NIAH-style writer warmup:
a probe head gives direct $\Mc\!\rightarrow$ answer supervision before
the LM-head path is allowed to dominate.
\end{enumerate}

\section{Memory Residuals Architecture}

Memory Residuals augments a pretrained causal LM with a recurrent state
$\Mc \in \mathbb{R}^{K \times d}$, where $K=128$ slots and $d$ is the
backbone hidden size.  The state is updated at session boundaries, not
per token.  The architecture has three stages: extraction from the
completed session, a judge/write step that updates $\Mc$, and a readout
that injects $\mt$ into the next session's residual stream.

\paragraph{Extract.}
Given a session representation $\Ct \in \mathbb{R}^{N_t \times d}$,
a learned query bank $\Min$ cross-attends to produce a candidate memory:
\begin{equation}
\Mnew =
\mathrm{softmax}\!\left(
  \frac{\Min W_Q^{\mathrm{ext}}(\Ct W_K^{\mathrm{ext}})^\top}{\sqrt d}
\right)\Ct W_V^{\mathrm{ext}} .
\end{equation}
The v15 recipe uses a contextual extraction source
(\texttt{hidden\_14}) and an input RMSNorm before extraction.  That norm
was introduced after gradient audits showed unnormalized hidden states
could drive extract-key gradients roughly $50\times$ too large.

\paragraph{Judge And Write.}
The candidate $\Mnew$ competes with the previous state through a
slot-attention writer:
\begin{equation}
P=[\Mc^{t-1}\parallel \Mnew], \qquad
\Mc^t =
\mathrm{Judge}_{\theta}(P).
\end{equation}
Earlier versions used a cross-attention judge and later versions use a
slot-attention writer with orthogonal query initialization and slot
position embeddings.  The important empirical point is that the writer
can settle into a content-blind fixed point: $\Mc$ becomes useful as a
generic ``memory is on'' prior while carrying little chain-specific
content.

\paragraph{Read.}
When scoring the next session, token representations $X$ attend to
$\Mc$:
\begin{equation}
\mt =
\mathrm{softmax}\!\left(
  \frac{XW_Q^{\mathrm{r}}(\Mc W_K^{\mathrm{r}})^\top}{\sqrt d}
\right)\Mc W_V^{\mathrm{r}} .
\end{equation}
The readout $\mt$ has the same shape as a residual-stream tensor and is
registered as an extra source in a Block-Attention Residual router.  A
learned routing mass $\amem$ determines how much the LM computation uses
the memory source at each routed sublayer.  This design is intended to
be additive on top of the backbone, but the experiments below show that
``additive'' is not enough: every non-memory route to the callback value
must be audited.

\section{Failure Modes At A Glance}

Table~\ref{tab:failure_modes} summarizes the six failure modes surfaced
by the campaign.  The first four can make memory look effective when it
is not, or hide a localized effect when it is real.  The fifth prevents
a false dismissal of the metric.  The sixth is the remaining
architecture/training problem after the shortcut paths are removed.

\begin{table}[t]
\centering
\caption{Failure modes and diagnostics.  ``Fix'' means the methodology
we adopt in this paper, not a claim that the architecture is solved.}
\label{tab:failure_modes}
\small
\begin{tabular}{p{0.18\linewidth}p{0.26\linewidth}p{0.25\linewidth}p{0.20\linewidth}}
\toprule
Name & Symptom & Diagnostic & Fix \\
\midrule
Token-averaged dilution &
Standard eval reports $\dnm \approx 0$ although callback tokens improve. &
Callback-token-only CE; compare score-token counts. &
Use \texttt{eval\_callback.py}. \\
\midrule
Joint-training invisibility &
Unfrozen backbone collapses $\elift$ to $\approx 0$. &
Freeze-vs-joint ablation; compare to empirical prior. &
Report memory-disabled and prior baselines. \\
\midrule
Category-cue leak &
``favorite COLOR'' reduces answer to a 32-way prior. &
$P(\mathrm{item}\mid\mathrm{category})$ CE = 2.91 nats; v15b CE = 2.97. &
Random IDs; remove category cue. \\
\midrule
Window leakage &
Evidence directly enters the LM attention window. &
Count evidence positions in callback windows. &
Use smaller callback window or forced evidence gap. \\
\midrule
Redaction uncertainty &
Zero lift could be dismissed as a leaky floor. &
Default/cross/skip/zero redaction agreement within 0.05 nats. &
Keep evidence-redacted floors and audit them. \\
\midrule
Writer-objective failure &
D5 removes priors but $\elift$ remains near zero. &
D5 run plus TTT-readout/probe audits. &
F2 NIAH writer warmup. \\
\bottomrule
\end{tabular}
\end{table}

\begin{figure}[t]
\centering
\fbox{\begin{minipage}{0.94\linewidth}
\small
\textbf{Five pathways from $\Mc$ or context to the LM head.}
\begin{enumerate}
\item Intended path: evidence session $\rightarrow \Ct \rightarrow$ writer $\rightarrow \Mc \rightarrow \mt \rightarrow$ LM head.
\item Content-blind memory prior: writer emits a similar $\Mc$ on every chain; readout shifts the LM toward the callback marginal.
\item Backbone template prior: unfrozen LM learns $P(\mathrm{item}\mid\mathrm{category})$ from callback tokens.
\item Window path: evidence remains in the local attention window, bypassing $\Mc$.
\item Measurement path: token-averaged eval or unaudited redaction changes the apparent contribution of memory.
\end{enumerate}
\end{minipage}}
\caption{Causal pathways that must be separated before a memory result
can be interpreted.  Only the first path is chain-specific recurrent
memory.}
\label{fig:pathways}
\end{figure}

\section{Failure Modes 1--2: Dilution And Joint-Training Invisibility}

\paragraph{Callback Benefit Is Localized.}
The first apparent contradiction was v14k.  The in-trainer and
callback-aware metrics showed a large callback-token benefit, while
standard chain evaluation made the same checkpoint look neutral or
worse.  The callback-aware post-hoc eval on 200 D4v2 validation chains
reports
\[
\CEmem=4.98,\quad
\CEnomem=6.42,\quad
\CEsh=5.04,\quad
\CEfloor=5.05.
\]
Thus $\dnm^{\mathrm{cb}}=+1.44$, $\dsh^{\mathrm{cb}}=+0.054$, and
$\elift=+0.071$.  The full-window standard eval reports $\dnm=-0.085$
on the same checkpoint.  These are not inconsistent: the callback mask
is a tiny part of the score window, and memory can add small noise on
non-callback tokens that dominates a full-window average.

\paragraph{Most Of The Apparent Benefit Survives Redaction.}
The same v14k number also reveals a second problem.  The evidence-
redacted floor still gives $\Delta_{\mathrm{nm-m,floor}}=+1.37$.  Therefore
$1.37/1.44 = 95.1\%$ of the apparent no-memory vs memory gain survives
when $\Mc$ is built without the answer.  The honest chain-specific
quantity is not $\dnm$ but
$\elift=\dnm-\Delta_{\mathrm{nm-m,floor}}$.

\paragraph{Joint Training Changes The Object Of Measurement.}
The v15 matrix makes the point sharper.  With a frozen backbone, the
memory path can change callback CE because the LM head cannot simply
fit the D4v2 closed-set marginal.  With an unfrozen backbone, the LM
path can learn the callback distribution directly.  v15b, the 0.6B
joint-trained run, reaches $\CEmem \approx 2.97$ nats and
$\elift \approx 0$.  This is not evidence that memory is unnecessary in
general; it is evidence that the benchmark allowed a non-memory path to
solve the proxy objective.

\begin{table}[t]
\centering
\caption{D4v2 ablation and audit results.  CE/prior rows are from the
A3 template-prior audit on 128 validation chains unless noted.  Lift
values mix in-trainer and post-hoc audits where indicated; missing
trajectory details are left as TODOs rather than interpolated.}
\label{tab:v15_ablation}
\small
\begin{tabular}{llrrrrp{0.24\linewidth}}
\toprule
Run & Regime & $\CEmem$ & Gap to 2.91 & $\dnm^{\mathrm{cb}}$ & $\elift$ & Reading \\
\midrule
v14k & 0.6B frozen & 4.98 & +2.07 & +1.44 & +0.071 & Large apparent callback gain; 95\% survives redaction. \\
v15a & 0.6B frozen & 5.15 & +2.24 & +1.33 & -0.043 & Reproduces $\dnm$; post-hoc floor is as strong as evidence. \\
v15b & 0.6B joint & 2.97 & +0.06 & +0.007 & -0.002 & Matches empirical prior; memory contribution vanishes. \\
v15c & 0.6B frozen, embed extract & \todo{NUMBER} & \todo{NUMBER} & \todo{NUMBER} & +0.005--+0.020 & Hidden-state extraction is better. \\
v15e & 1.7B frozen & 4.76 & +1.85 & \todo{NUMBER} & +0.066 / -0.09 & Strong $\dnm$, weak or unstable evidence-specific lift. \\
v15f & 1.7B joint & 3.21 & +0.30 & \todo{NUMBER} & $\approx$+0.04 & Converging toward the same prior-matching collapse. \\
\bottomrule
\end{tabular}
\end{table}

\section{Failure Mode 3: The Category-Cue Corpus Leak}

D4v2 was meant to test whether $\Mc$ stores a persona binding.  Its
callback template inadvertently answers a different question:
can the model learn a category-conditional closed-set prior?  The
callback repeats the category in the user question and immediately
before the answer:
\begin{quote}
\small
User: Quick question, what was my favorite COLOR again?\\
Assistant: Your favorite COLOR is RED.
\end{quote}
Given the category, the answer space is roughly 32 items.  Audit A3
builds a Laplace-smoothed empirical $P(\mathrm{item}\mid\mathrm{category})$
from the 5000-chain training set and scores 128 validation callbacks.
The resulting prior CE is $2.9087$ nats per callback-mask token.  The
joint 0.6B v15b checkpoint scores $2.9655$ nats at best and $2.9394$
nats at final, within $0.06$ and $0.03$ nats of the prior.  That is the
smoking gun: the unfrozen backbone has learned the template prior, not
retrieval.

Frozen-base controls rule out the opposite explanation that Qwen3
already knew this prior.  Without memory or fine-tuning, callback-only
CE is $7.69$ nats for Qwen3-0.6B and $8.76$ nats for Qwen3-1.7B on the
same masked tokens.  The prior is learned during the chain training
regime.

\begin{table}[t]
\centering
\caption{D4v2 versus D5 corpus comparison.  D5 removes the category cue
and uses per-chain random alphanumeric codes; observed v16a results
show that removing the prior exposes, but does not solve, the writer
objective problem.}
\label{tab:corpus_compare}
\small
\begin{tabular}{lrrrp{0.28\linewidth}}
\toprule
Corpus & Prior / floor CE & Base CE & Observed $\elift$ & Interpretation \\
\midrule
D4v2 closed-set persona & 2.91 & 7.69 / 8.76 & v14k +0.071; v15b $\approx$0 & Category cue makes prior-fitting sufficient under joint training. \\
D5 random codes & $\sim$10 nats per ID-character; first observed callback CE $\sim$7 & \todo{NUMBER} & v16a -0.002 at step 100; later +0.011 & Prior removed; writer still does not encode chain-specific binding. \\
\bottomrule
\end{tabular}
\end{table}

\section{Failure Mode 4: Random IDs Are Not A Cure}

D5 implements the obvious fix.  It replaces closed-set persona items
with per-chain random alphanumeric codes and removes the category cue.
The evidence session acknowledges the code without echoing it in the
assistant response, and \texttt{--mask\_evidence\_session\_loss} zeros
LM loss over evidence sessions.  In principle, the callback answer has
no useful non-memory path: the code is random, the callback does not
identify a category, and the LM head is not directly trained on the
evidence answer span.

The first D5 result is therefore important precisely because it is
negative.  v16a, a frozen 0.6B D5 run, reports
$\mathrm{evidence\_lift}=-0.0024$ at step 100; the v17 proposal packet
summarizes the later state as $\mathrm{evidence\_lift}=+0.011$ with
$\dnm^{\mathrm{cb}}=-0.076$, InfoNCE loss $1.375 \approx \log 4$, D3-MC
pair/self $0.012$, and $\amem \approx 0.01$.  Removing the prior made
the evaluation honest, but it did not create a training signal that
forces the writer to bind the code into $\Mc$.

This is a different failure from the D4v2 category leak.  On D4v2 the
LM path can win by learning a closed-set marginal.  On D5 the marginal
is unhelpful, but the writer still receives its gradient through a
long, weak path: callback NLL $\rightarrow$ LM head $\rightarrow$
depth-router $\rightarrow$ readout $\rightarrow \Mc \rightarrow$ writer.
Near random initialization, wrong codes and uninformative memories look
similarly bad to the LM loss.  The writer has no direct extractive
objective saying, ``the code in this evidence session must be present in
$\Mc$.''  That is the motivation for F2.

\section{Failure Mode 5: Leak Audits As Methodology}

An honest negative result requires proving that the diagnostic itself is
not broken.  Audit A2 traces both the in-trainer floor and standalone
callback floor.  The in-trainer floor builds $\Mc$ fresh from a single
non-evidence session in the same chain; the standalone floor walks the
prefix while replacing evidence sessions with distractors.  Neither
path carries a cached rolling state across evaluation passes.

The adversarial redaction test is stronger than the code trace.  On
v15a/best, the default redaction gives $\CEfloor=5.1235$; cross-chain
redaction gives $5.0911$; skip-evidence gives $5.0857$; zeroed-input
redaction gives $5.0896$.  The spread is $0.038$ nats, inside the
pre-declared $0.05$-nat band, and the direction is opposite a leak:
the default floor is slightly worse, not better, than cross/zero.  On
v15b/best, all floor variants lie within $0.008$ nats.  Thus
$\elift \approx 0$ is not a redaction artifact.

Audit A1 adds a separate lesson.  At D4v2's v15 default
\texttt{window\_k=4}, $58.5\%$ of train chains and $59.6\%$ of
validation chains have at least one evidence session inside the LM
attention window at callback scoring time.  At \texttt{window\_k=3},
the leak is still $41.4\%$ / $43.4\%$.  Only \texttt{window\_k=1}
eliminates direct evidence visibility.  This makes callback-window
counts a required part of future corpus reports.

\begin{figure}[t]
\centering
\fbox{\begin{minipage}{0.94\linewidth}
\small
\centering
\begin{tabular}{lrrrr}
\toprule
Run & Corpus & Scale/regime & $\dnm^{\mathrm{cb}}$ & $\elift$ \\
\midrule
v14k & D4v2 & 0.6B frozen & +1.44 & +0.071 \\
v15a & D4v2 & 0.6B frozen & +1.33 & -0.043 (post-hoc floor) \\
v15b & D4v2 & 0.6B joint & +0.007 & -0.002 \\
v15e & D4v2 & 1.7B frozen & +2.03 (trainer) & +0.066 / -0.09 \\
v15f & D4v2 & 1.7B joint & \todo{NUMBER} & $\approx$+0.04 \\
v16a & D5 & 0.6B frozen & -0.076 & +0.011 \\
v16b & D5 & 1.7B frozen & \todo{NUMBER} & \todo{NUMBER} \\
\bottomrule
\end{tabular}
\end{minipage}}
\caption{Trajectory summary for the main phase-aligned callback
metric.  The underlying logs are heterogeneous; this figure therefore
shows audited checkpoints rather than an interpolated curve.}
\label{fig:trajectory}
\end{figure}

\section{Prescription: F2 NIAH Writer Warmup}

The F2 prescription is to add a writer-side extractive objective before
the LM head is allowed to use memory.  The proposed v17a/v17e warmup
uses the D5 random-code corpus and a probe head that reads $\Mc$ and
predicts the code.  The probe is discarded at evaluation time; it is a
training scaffold, not an inference path.  If the final LM-head
callback evaluation improves after the probe is removed, the result
means the writer has learned to store content that the normal readout
can use.

The warmup tests two hypotheses cleanly.  If probe CE drops well below
the D5 marginal but downstream $\elift$ remains near zero, then the
writer can encode the answer and the bottleneck is readout/router use.
If probe CE stays at the marginal, the writer cannot encode the binding
even under direct supervision and the architecture itself needs a
writer-side change.  Either outcome is more informative than another
LM-only training run.

\begin{table}[t]
\centering
\caption{D5 TTT-readout audit placeholder.  The parent run will fill
these values when \texttt{tools/d5\_ttt\_readout.py} completes.  All
unknown numbers are deliberately left as \texttt{\textbackslash todo}.}
\label{tab:ttt_readout}
\small
\begin{tabular}{lrrrrp{0.25\linewidth}}
\toprule
Checkpoint & Probe params trained & Base CE & TTT CE & TTT $\elift$ & Verdict \\
\midrule
v15a/best & $W_Q,W_K,W_V$ only & \todo{NUMBER} & \todo{NUMBER} & \todo{NUMBER} & \todo{writer vs readout bottleneck} \\
v15e/best & $W_Q,W_K,W_V$ only & \todo{NUMBER} & \todo{NUMBER} & \todo{NUMBER} & \todo{writer vs readout bottleneck} \\
\bottomrule
\end{tabular}
\end{table}

\begin{figure}[t]
\centering
\fbox{\begin{minipage}{0.94\linewidth}
\small
\centering
\begin{tabular}{lrrrr}
\toprule
Run & Scale & Warmup/probe CE & $\dnm^{\mathrm{cb}}$ & $\elift$ \\
\midrule
v17a & 0.6B & \todo{NUMBER} & \todo{NUMBER} & \todo{NUMBER} \\
v17e & 1.7B & \todo{NUMBER} & \todo{NUMBER} & \todo{NUMBER} \\
\bottomrule
\end{tabular}
\end{minipage}}
\caption{Placeholder for v17 F2 trajectories.  Positive $\elift$ after
probe removal would close the loop: D5 removes the shortcut and F2
supplies the missing writer gradient.}
\label{fig:v17}
\end{figure}

\section{Related Work}

Our findings align with a broad lesson from long-context and memory
evaluation: shortcuts are the default when the benchmark permits them.
NoLiMa removes literal overlap between query and evidence and shows
that ordinary long-context retrieval performance can be template-driven
rather than robust memory use~\cite{nolima2025}.  RULER argues that
vanilla needle-in-a-haystack is saturated and adds UUIDs, multiple
needles, variable tracking, aggregation, and hard distractors to expose
weaker retrieval behavior~\cite{ruler2024}.  Lost in the Middle shows
position and recency/primacy effects even when context length is
nominally large~\cite{lostmiddle2024}.  LongBench includes no-context
checks to expose parametric priors~\cite{longbench2024}; our
callback-only and empirical-prior audits are the memory analogue.
LongMemEval separates indexing, retrieval, and reading, motivating our
separation of write, readout, and decode diagnostics~\cite{longmemeval2025}.

Memory architectures face the same credit-assignment issue.  Recurrent
Memory Transformer, Associative RMT, AutoCompressors, Memorizing
Transformers, TRIME, Block-Recurrent Transformers, and classical
differentiable memory systems all introduce pathways that can, in
principle, carry long-range information~\cite{rmt2022,armt2024,
autocompressors2023,memorizing2022,trime2022,blockrecurrent2022,
ntm2014}.  Our results emphasize that adding a pathway is not the same
as proving the model used it.  The evaluation must include memory-off,
memory-shuffled, evidence-redacted, prior-only, and window-leakage
controls.

\section{Conclusion And Open Questions}

The campaign gives a conservative answer to ``when does recurrent memory
help?''  It helps only when the task and evaluator make the non-memory
paths uninformative and when the writer receives a gradient that can
actually encode the binding.  D4v2 failed the first condition: category
cues, window leakage, and token-averaged eval made the memory result
ambiguous.  D5 satisfies the first condition better and exposes the
second: LM-only callback loss is too weak to train the writer out of a
content-blind fixed point.

The next decisive experiment is F2.  A positive v17a/v17e result would
show that direct writer supervision plus shortcut-free data yields
chain-specific memory under the normal LM-head readout.  A negative
result would be equally useful: it would localize the failure to the
readout/router or the memory representation itself.  Either way, the
methodological bar is now clear: report the shortcuts before reporting
the memory gain.

\section*{Reproducibility Statement}

All numbers in the main text are drawn from the run ledger and audits in
\texttt{results/exp2\_chain\_recipe/} or from callback-aware JSON
outputs in \texttt{results/eval\_v14v15/}.  Missing values are marked
with \todo{NUMBER}.  The relevant tools are
\texttt{tools/eval\_callback.py}, \texttt{tools/audit\_a1\_window\_leakage.py},
\texttt{tools/audit\_a2\_redaction.py}, \texttt{tools/audit\_a3\_template\_prior.py},
\texttt{tools/build\_synthetic\_random\_codes.py}, and the planned
\texttt{tools/d5\_ttt\_readout.py}.

\begin{thebibliography}{99}

\bibitem{nolima2025}
Modarressi et al.
\newblock NoLiMa: Long-Context Evaluation Beyond Literal Matching.
\newblock arXiv:2502.05167.

\bibitem{ruler2024}
Hsieh et al.
\newblock RULER: What's the Real Context Size of Your Long-Context Language Models?
\newblock arXiv:2404.06654.

\bibitem{lostmiddle2024}
Liu et al.
\newblock Lost in the Middle: How Language Models Use Long Contexts.
\newblock arXiv:2307.03172.

\bibitem{longbench2024}
Bai et al.
\newblock LongBench: A Bilingual, Multitask Benchmark for Long Context Understanding.
\newblock arXiv:2308.14508.

\bibitem{babilong2024}
Kuratov et al.
\newblock BABILong: Testing the Limits of LLMs with Long Context Reasoning-in-a-Haystack.
\newblock arXiv:2406.10149.

\bibitem{longmemeval2025}
Yu et al.
\newblock LongMemEval: Benchmarking Chat Assistants on Long-Term Interactive Memory.
\newblock arXiv:2410.10813.

\bibitem{shortcut2020}
Geirhos et al.
\newblock Shortcut Learning in Deep Neural Networks.
\newblock arXiv:2004.07780.

\bibitem{memorizing2022}
Wu et al.
\newblock Memorizing Transformers.
\newblock arXiv:2203.08913.

\bibitem{rmt2022}
Bulatov, Kuratov, and Burtsev.
\newblock Recurrent Memory Transformer.
\newblock arXiv:2207.06881.

\bibitem{armt2024}
Rodkin et al.
\newblock Associative Recurrent Memory Transformer.
\newblock ICML 2024 Next Generation of Sequence Modeling Architectures workshop.

\bibitem{autocompressors2023}
Chevalier et al.
\newblock Adapting Language Models to Compress Contexts.
\newblock arXiv:2305.14788.

\bibitem{trime2022}
Zhong, Lei, and Chen.
\newblock Training Language Models with Memory Augmentation.
\newblock arXiv:2205.12674.

\bibitem{blockrecurrent2022}
Hutchins et al.
\newblock Block-Recurrent Transformers.
\newblock NeurIPS 2022.

\bibitem{ntm2014}
Graves, Wayne, and Danihelka.
\newblock Neural Turing Machines.
\newblock arXiv:1410.5401.

\bibitem{zeroscrolls2023}
Shaham et al.
\newblock ZeroSCROLLS: A Zero-Shot Benchmark for Long Text Understanding.
\newblock arXiv:2305.14196.

\bibitem{ragbaseline2025}
Laitenberger et al.
\newblock Stronger Baselines for Retrieval-Augmented Generation with Long-Context Language Models.
\newblock arXiv:2506.03989.

\end{thebibliography}

\appendix

\section{Audit Details}

\subsection{Window Leakage Counts}
For D4v2, evidence positions are sampled from body positions $0,\ldots,8$
and the callback is at position 9.  Audit A1 counts whether an evidence
session falls in the callback scoring window.  At \texttt{window\_k=4},
2926/5000 train chains and 298/500 validation chains leak at least one
evidence session.  At \texttt{window\_k=3}, the corresponding counts are
2068/5000 and 217/500.  At \texttt{window\_k=2}, they are 1088/5000 and
114/500.  Only \texttt{window\_k=1} eliminates this direct path.

\subsection{Redaction Variants}
Audit A2 evaluates five floor variants: original evidence (match),
default distractor-pool redaction, cross-chain non-evidence redaction,
skip-evidence, and zeroed-token evidence.  On v15a/best, the
redaction-style spread is 0.038 nats; on v15b/best it is 0.0066 nats.
The default floor is slightly worse than cross/zero, so the default is
not hiding evidence in the dangerous direction.

\subsection{Template Prior Construction}
Audit A3 parses $(\mathrm{category},\mathrm{item})$ from D4v2 chain
names and builds an add-half-smoothed empirical prior over each
category.  The train counts range from 593 to 652 chains per category.
The validation prior CE over the first 128 chains is $2.9087$ nats per
callback-mask token.  The per-item mean is essentially $\log 32$; the
per-token mean is lower because 40/128 validation items are multi-token.

\section{Callback-Aware Metric Definition}

Let $S_c$ be the callback session of chain $c$ and $A_c$ its annotated
answer-token mask.  For a memory construction $r$, define
\[
\mathrm{CE}_r^{\mathrm{cb}}(c)=
\frac{1}{|A_c|}\sum_{i\in A_c}
  -\log p_\theta(x_i \mid x_{<i}, M_{\mathrm{c},r}(c)).
\]
The reported callback deltas are
\[
\dnm^{\mathrm{cb}}=
\mathrm{CE}_{\mathrm{nomem}}^{\mathrm{cb}}-
\mathrm{CE}_{\mathrm{mem}}^{\mathrm{cb}},\qquad
\dsh^{\mathrm{cb}}=
\mathrm{CE}_{\mathrm{shuffle}}^{\mathrm{cb}}-
\mathrm{CE}_{\mathrm{mem}}^{\mathrm{cb}},
\]
and
\[
\elift =
\dnm^{\mathrm{cb}} -
\left(
\mathrm{CE}_{\mathrm{nomem,floor}}^{\mathrm{cb}}-
\mathrm{CE}_{\mathrm{mem,floor}}^{\mathrm{cb}}
\right).
\]
Positive $\elift$ means evidence-bearing $\Mc$ helps more than
evidence-redacted $\Mc$.

\section{Additional Tables}

\begin{table}[h]
\centering
\caption{Base callback CE without memory or fine-tuning on D4v2
callback-mask tokens.}
\label{tab:base_ce}
\small
\begin{tabular}{lr}
\toprule
Backbone & CE on callback tokens \\
\midrule
Qwen3-0.6B base & 7.686 \\
Qwen3-1.7B base & 8.764 \\
Uniform over 256 closed-set items & 5.545 \\
Uniform within cued 32-item category & 3.466 \\
Empirical per-token template prior & 2.909 \\
\bottomrule
\end{tabular}
\end{table}

\begin{table}[h]
\centering
\caption{A2 redaction audit summary.}
\label{tab:redaction_summary}
\small
\begin{tabular}{lrrrr}
\toprule
Checkpoint & Match $\dnm$ & Default $\elift$ & Cross $\elift$ & Redaction spread \\
\midrule
v15a/best & +1.3377 & -0.0310 & -0.0634 & 0.038 \\
v15b/best & +0.0074 & -0.0021 & -0.0078 & 0.0066 \\
\bottomrule
\end{tabular}
\end{table}

\section{Limitations}

The paper is a first draft under an active experiment campaign.  Several
v17 and TTT-readout numbers are intentionally left as placeholders.  The
D5 corpus removes the most obvious D4v2 prior but still uses synthetic
templates; future natural-memory benchmarks such as LoCoMo and MSC
should be reported with the same no-memory, prior, redaction, and
shuffle controls.  Finally, Memory Residuals itself is only one
architecture.  The central claim of this paper is methodological:
without these controls, any recurrent-memory architecture can appear to
work through hidden pathways.

\end{document}
